\documentclass[11pt,a4paper]{article}
\usepackage[margin=2.5cm]{geometry}
\usepackage[T1]{fontenc}
\usepackage[utf8]{inputenc}
\usepackage{lmodern}
\usepackage{amsmath,amssymb,amsthm,mathtools,physics,bm}
\usepackage{microtype}
\title{Lösungen}
\date{}
\begin{document}
\maketitle

\section*{Aufgabe 1 [v4]}
\subsection*{(1.1) [v4]}
\paragraph{Aufgabentext.}
Natural units, $\hbar=c=1$ : how much is 1 kg in GeV and $1 s$ in $\mathrm{GeV}^{-1}$ ? Use your result to express Newton's gravitational constant $G_N=6.67 \times 10^{-11} \mathrm{~m}^3 \mathrm{~kg}^{-1} \mathrm{~s}^{-2}$ in natural units, and give the value of the Planck mass $M_{P l}=G_N^{-1 / 2}$.

\paragraph{Lösung.}
Zunächst konvertieren wir 1 kg in GeV. In natürlichen Einheiten gilt $\hbar = c = 1$, was bedeutet, dass Energie, Masse und Impuls dieselben Einheiten haben. Die Umrechnung von kg in GeV erfolgt über die Ruheenergie eines Kilogramms, $E = mc^2$. Da $c = 1$ in natürlichen Einheiten, ist $c^2$ ebenfalls 1, und wir verwenden die Umrechnung von Joule in GeV:

\[
1 \, \text{kg} = \frac{1 \, \text{kg} \cdot (3 \times 10^8 \, \text{m/s})^2}{1.60218 \times 10^{-10} \, \text{J/GeV}}
\]

Da $1 \, \text{J} = 1 \, \text{kg} \cdot \text{m}^2/\text{s}^2$, erhalten wir:

\[
1 \, \text{kg} \approx 5.61 \times 10^{26} \, \text{GeV}
\]

Nun konvertieren wir 1 s in $\mathrm{GeV}^{-1}$. In natürlichen Einheiten, wo $c = 1$, ist die Lichtgeschwindigkeit gleich 1, was bedeutet, dass Zeit und Länge dieselben Einheiten haben. Die Umrechnung erfolgt über die Beziehung zwischen Energie und Zeit, die durch die Planck-Konstante $\hbar$ vermittelt wird. Da $\hbar = 6.582 \times 10^{-25} \, \text{GeV s}$, ergibt sich:

\[
1 \, \text{s} = \frac{1}{6.582 \times 10^{-25} \, \text{GeV}}
\]

Jetzt berechnen wir die Gravitationskonstante $G_N$ in natürlichen Einheiten. In SI-Einheiten hat $G_N$ die Dimension $\mathrm{m}^3 \, \mathrm{kg}^{-1} \, \mathrm{s}^{-2}$. Um dies in natürlichen Einheiten umzurechnen, verwenden wir:

\[
G_N = 6.67 \times 10^{-11} \, \mathrm{m}^3 \, \mathrm{kg}^{-1} \, \mathrm{s}^{-2} \times \left( \frac{6.582 \times 10^{-25} \, \text{GeV}}{1} \right)^2 \times \left( \frac{1}{5.61 \times 10^{26} \, \text{GeV}} \right)
\]

Dies ergibt:

\[
G_N \approx 6.71 \times 10^{-39} \, \text{GeV}^{-2}
\]

Schließlich berechnen wir die Planck-Masse $M_{Pl}$:

\[
M_{Pl} = G_N^{-1/2} \approx (6.71 \times 10^{-39} \, \text{GeV}^{-2})^{-1/2} \approx 1.22 \times 10^{19} \, \text{GeV}
\]

\section*{Aufgabe 2 [v4]}
\subsection*{(2.1) [v4]}
\paragraph{Aufgabentext.}
Starting from the definition of the Hamiltonian

$$
H(x, p) \equiv p \dot{x}-L(x, \dot{x})
$$

derive Hamilton's equations

$$
\frac{\partial H}{\partial x}=-\dot{p}, \quad \frac{\partial H}{\partial p}=\dot{x}
$$

[Hint: Be careful about what the independent variables of a function are.]

\paragraph{Lösung.}
To derive Hamilton's equations from the given Hamiltonian, we start by considering the Hamiltonian function:

$$
H(x, p) = p \dot{x} - L(x, \dot{x}),
$$

where \( \dot{x} \) is expressed as a function of \( x \) and \( p \) through the Legendre transformation. The relationship between \( \dot{x} \) and \( p \) is given by:

$$
p = \frac{\partial L}{\partial \dot{x}}.
$$

To find Hamilton's equations, we need to compute the partial derivatives of \( H \) with respect to \( x \) and \( p \).

**1. Derivative of \( H \) with respect to \( x \):**

Using the chain rule, we have:

$$
\frac{\partial H}{\partial x} = \frac{\partial}{\partial x} (p \dot{x} - L) = p \frac{\partial \dot{x}}{\partial x} - \frac{\partial L}{\partial x}.
$$

In the context of the Hamiltonian formulation, \( \dot{x} \) is considered a function of \( p \) and \( x \), but when taking the partial derivative of \( H \) with respect to \( x \), \( \dot{x} \) is treated as independent of \( x \). This is because \( \dot{x} \) is determined by the Legendre transformation and is not an explicit function of \( x \) in the partial derivative. Thus, the term involving \( \frac{\partial \dot{x}}{\partial x} \) vanishes, leading to:

$$
\frac{\partial H}{\partial x} = -\frac{\partial L}{\partial x}.
$$

The Euler-Lagrange equation gives:

$$
\frac{d}{dt} \left( \frac{\partial L}{\partial \dot{x}} \right) = \frac{\partial L}{\partial x}.
$$

Substituting \( \frac{\partial L}{\partial \dot{x}} = p \), we obtain:

$$
\frac{d}{dt}(p) = \frac{\partial L}{\partial x}.
$$

Thus, the correct expression for the time derivative of \( p \) is:

$$
\dot{p} = -\frac{\partial H}{\partial x}.
$$

**2. Derivative of \( H \) with respect to \( p \):**

For the partial derivative with respect to \( p \), we have:

$$
\frac{\partial H}{\partial p} = \frac{\partial}{\partial p} (p \dot{x} - L) = \dot{x},
$$

since \( \dot{x} \) is treated as a function of \( p \) and \( x \).

Therefore, Hamilton's equations are:

$$
\dot{x} = \frac{\partial H}{\partial p}, \quad \dot{p} = -\frac{\partial H}{\partial x}.
$$

These equations describe the time evolution of the system in terms of the Hamiltonian function.

\section*{Aufgabe 3 [v4]}
\subsection*{(3.1) [v4]}
\paragraph{Aufgabentext.}
Using the Schrödinger equation for the wavefunction $\Psi(\mathbf{x}, t)$

$$
\left\{-\frac{\hbar^2 \nabla^2}{2 m}+V(\mathbf{x})\right\} \Psi(\mathbf{x}, t)=i \hbar \frac{\partial}{\partial t} \Psi(\mathbf{x}, t)
$$

show that the probability density $\rho=\Psi^* \Psi$ satisfies the continuity equation

$$
\frac{\partial}{\partial t} \rho+\nabla \cdot \mathbf{j}=0
$$

where

$$
\mathbf{j}=\frac{\hbar}{2 i m}\left\{\Psi^* \nabla \Psi-\left(\nabla \Psi^*\right) \Psi\right\} .
$$

[Hint: Consider $\Psi^* \times$ (Schr.Eq.) $-\Psi \times$ (Schr.Eq.)*]

\paragraph{Lösung.}
Um zu zeigen, dass die Wahrscheinlichkeitsdichte \(\rho = \Psi^* \Psi\) die Kontinuitätsgleichung erfüllt, beginnen wir mit der Schrödinger-Gleichung und ihrer komplex konjugierten Form:

\[
\left\{-\frac{\hbar^2 \nabla^2}{2 m} + V(\mathbf{x})\right\} \Psi = i \hbar \frac{\partial \Psi}{\partial t}
\]

\[
\left\{-\frac{\hbar^2 \nabla^2}{2 m} + V(\mathbf{x})\right\} \Psi^* = -i \hbar \frac{\partial \Psi^*}{\partial t}
\]

Multiplizieren wir die Schrödinger-Gleichung mit \(\Psi^*\) und die konjugierte Schrödinger-Gleichung mit \(\Psi\):

\[
\Psi^* \left\{-\frac{\hbar^2 \nabla^2}{2 m} \Psi + V(\mathbf{x}) \Psi\right\} = \Psi^* i \hbar \frac{\partial \Psi}{\partial t}
\]

\[
\Psi \left\{-\frac{\hbar^2 \nabla^2}{2 m} \Psi^* + V(\mathbf{x}) \Psi^*\right\} = -\Psi i \hbar \frac{\partial \Psi^*}{\partial t}
\]

Subtrahieren wir die zweite Gleichung von der ersten:

\[
\Psi^* \left(-\frac{\hbar^2 \nabla^2}{2 m} \Psi\right) - \Psi \left(-\frac{\hbar^2 \nabla^2}{2 m} \Psi^*\right) = \Psi^* i \hbar \frac{\partial \Psi}{\partial t} + \Psi i \hbar \frac{\partial \Psi^*}{\partial t}
\]

Vereinfachen wir die linke Seite:

\[
-\frac{\hbar^2}{2 m} \left(\Psi^* \nabla^2 \Psi - \Psi \nabla^2 \Psi^*\right) = i \hbar \left(\Psi^* \frac{\partial \Psi}{\partial t} + \Psi \frac{\partial \Psi^*}{\partial t}\right)
\]

Die linke Seite kann umgeschrieben werden als:

\[
\nabla \cdot \left(\frac{\hbar}{2 i m} \left(\Psi^* \nabla \Psi - \Psi \nabla \Psi^*\right)\right)
\]

Diese Umformung zur Divergenzform ist korrekt, da der Ausdruck \(\Psi^* \nabla^2 \Psi - \Psi \nabla^2 \Psi^*\) durch Anwendung der Produktregel umgeschrieben werden kann. Die Produktregel besagt, dass:

\[
\nabla \cdot (\Psi^* \nabla \Psi) = \Psi^* \nabla^2 \Psi + \nabla \Psi^* \cdot \nabla \Psi
\]

\[
\nabla \cdot (\Psi \nabla \Psi^*) = \Psi \nabla^2 \Psi^* + \nabla \Psi \cdot \nabla \Psi^*
\]

Subtrahiert man diese, erhält man:

\[
\nabla \cdot (\Psi^* \nabla \Psi - \Psi \nabla \Psi^*) = \Psi^* \nabla^2 \Psi - \Psi \nabla^2 \Psi^* + \nabla \Psi^* \cdot \nabla \Psi - \nabla \Psi \cdot \nabla \Psi^*
\]

Da die Terme \(\nabla \Psi^* \cdot \nabla \Psi\) und \(\nabla \Psi \cdot \nabla \Psi^*\) sich gegenseitig aufheben, bleibt:

\[
\nabla \cdot (\Psi^* \nabla \Psi - \Psi \nabla \Psi^*) = \Psi^* \nabla^2 \Psi - \Psi \nabla^2 \Psi^*
\]

Setzen wir dies in die Kontinuitätsgleichung ein:

\[
\frac{\partial}{\partial t} (\Psi^* \Psi) + \nabla \cdot \left(\frac{\hbar}{2 i m} \left(\Psi^* \nabla \Psi - \Psi \nabla \Psi^*\right)\right) = 0
\]

Dies zeigt, dass die Wahrscheinlichkeitsdichte \(\rho = \Psi^* \Psi\) die Kontinuitätsgleichung erfüllt, wobei der Stromdichtevektor \(\mathbf{j}\) gegeben ist durch:

\[
\mathbf{j} = \frac{\hbar}{2 i m} \left(\Psi^* \nabla \Psi - \Psi \nabla \Psi^*\right)
\]

\section*{Aufgabe 4 [v4]}
\subsection*{(4.1) [v4]}
\paragraph{Aufgabentext.}
Let $\hat{O}$ be a time-independent operator in the Schrödinger picture, and $\hat{H}$ the timeindependent Hamiltonian of the system. Starting from the definition of a Heisenberg operator,

$$
\hat{O}_H(t)=e^{\frac{i}{\hbar} \hat{H}\left(t-t_0\right)} \hat{O} e^{-\frac{i}{\hbar} \hat{H}\left(t-t_0\right)}
$$

derive the Heisenberg equation of motion

$$
i \hbar \frac{d \hat{O}_H}{d t}=\left[\hat{O}_H, \hat{H}_H\right]
$$

\paragraph{Lösung.}
Um die Heisenberg-Gleichung der Bewegung herzuleiten, beginnen wir mit der Definition des Heisenberg-Operators:

\[
\hat{O}_H(t) = e^{\frac{i}{\hbar} \hat{H}(t-t_0)} \hat{O} e^{-\frac{i}{\hbar} \hat{H}(t-t_0)}
\]

Wir berechnen die zeitliche Ableitung von \(\hat{O}_H(t)\):

\[
\frac{d \hat{O}_H}{d t} = \frac{d}{d t} \left( e^{\frac{i}{\hbar} \hat{H}(t-t_0)} \hat{O} e^{-\frac{i}{\hbar} \hat{H}(t-t_0)} \right)
\]

Verwenden der Produktregel ergibt:

\[
\frac{d \hat{O}_H}{d t} = \left( \frac{d}{d t} e^{\frac{i}{\hbar} \hat{H}(t-t_0)} \right) \hat{O} e^{-\frac{i}{\hbar} \hat{H}(t-t_0)} + e^{\frac{i}{\hbar} \hat{H}(t-t_0)} \hat{O} \left( \frac{d}{d t} e^{-\frac{i}{\hbar} \hat{H}(t-t_0)} \right)
\]

Die Ableitungen der Exponentialoperatoren sind:

\[
\frac{d}{d t} e^{\frac{i}{\hbar} \hat{H}(t-t_0)} = \frac{i}{\hbar} \hat{H} e^{\frac{i}{\hbar} \hat{H}(t-t_0)}
\]

\[
\frac{d}{d t} e^{-\frac{i}{\hbar} \hat{H}(t-t_0)} = -\frac{i}{\hbar} \hat{H} e^{-\frac{i}{\hbar} \hat{H}(t-t_0)}
\]

Einsetzen dieser Ausdrücke in die Produktregel ergibt:

\[
\frac{d \hat{O}_H}{d t} = \frac{i}{\hbar} \hat{H} e^{\frac{i}{\hbar} \hat{H}(t-t_0)} \hat{O} e^{-\frac{i}{\hbar} \hat{H}(t-t_0)} - e^{\frac{i}{\hbar} \hat{H}(t-t_0)} \hat{O} \frac{i}{\hbar} \hat{H} e^{-\frac{i}{\hbar} \hat{H}(t-t_0)}
\]

Dies kann umgeschrieben werden als:

\[
\frac{d \hat{O}_H}{d t} = \frac{i}{\hbar} \left( \hat{H} \hat{O}_H(t) - \hat{O}_H(t) \hat{H} \right)
\]

Dies ist der Kommutator:

\[
\frac{d \hat{O}_H}{d t} = \frac{i}{\hbar} [\hat{H}, \hat{O}_H(t)]
\]

Multiplizieren mit \(i \hbar\) ergibt die Heisenberg-Gleichung der Bewegung:

\[
i \hbar \frac{d \hat{O}_H}{d t} = [\hat{O}_H, \hat{H}]
\]

Diese Herleitung gilt unter der Annahme, dass der Hamiltonoperator \(\hat{H}\) zeitunabhängig ist. Falls der Hamiltonoperator zeitabhängig ist, muss die Herleitung angepasst werden. In diesem Fall ist der Hamiltonoperator selbst eine Funktion der Zeit, \(\hat{H}(t)\), und die Ableitung der Exponentialoperatoren enthält zusätzliche Terme:

\[
\frac{d}{d t} e^{\frac{i}{\hbar} \int_{t_0}^{t} \hat{H}(t') dt'} = \frac{i}{\hbar} \hat{H}(t) e^{\frac{i}{\hbar} \int_{t_0}^{t} \hat{H}(t') dt'}
\]

\[
\frac{d}{d t} e^{-\frac{i}{\hbar} \int_{t_0}^{t} \hat{H}(t') dt'} = -\frac{i}{\hbar} \hat{H}(t) e^{-\frac{i}{\hbar} \int_{t_0}^{t} \hat{H}(t') dt'}
\]

Diese zusätzlichen Terme müssen in die Herleitung einbezogen werden, was zu einer modifizierten Form der Heisenberg-Gleichung führt. Im Heisenberg-Bild bleibt der Hamiltonoperator \(\hat{H}_H = \hat{H}(t)\) konstant, wenn er zeitunabhängig ist, andernfalls ist er explizit zeitabhängig. Die modifizierte Heisenberg-Gleichung für einen zeitabhängigen Hamiltonoperator lautet:

\[
i \hbar \frac{d \hat{O}_H}{d t} = [\hat{O}_H, \hat{H}(t)] + i \hbar \left(\frac{\partial \hat{O}_H}{\partial t}\right)
\]

Hierbei berücksichtigt der zusätzliche Term \(\frac{\partial \hat{O}_H}{\partial t}\) die explizite Zeitabhängigkeit des Operators \(\hat{O}_H\).

\end{document}\n